\section{Training pipeline}
\label{sec:method-training}

The structured M\textsuperscript{3}N predictor of
Section~\ref{sec:method-architecture} is trained through a
unified pipeline: one \texttt{Trainer} class, two losses, one
fit loop. This section describes the two training paths
selectable via YAML
(Section~\ref{subsec:method-training-paths}), the optimizer
and learning-rate schedule
(Section~\ref{subsec:method-training-optimizer}), the
per-example dual-variable bank that LP-M3N requires
(Section~\ref{subsec:method-training-phi}), and the
validation-based hyperparameter-selection protocol
(Section~\ref{subsec:method-training-hyperparams}).

\subsection{Training paths}
\label{subsec:method-training-paths}

Two structured losses are selectable via the YAML
\texttt{training.loss} field; the unified \texttt{Trainer}
dispatches between them based on this single string.

\paragraph{Structured hinge with Viterbi loss-augmented
inference (\texttt{m3n\_hinge}).}
The training loss is the margin-rescaling form
$\Delta(x, y, w)$ of Section~\ref{subsec:max-margin}, with the
loss-augmented oracle solved exactly on chains by Viterbi. The
corresponding implementation is the
\texttt{structured\_hinge\_loss} module paired with
\texttt{loss\_augmented\_viterbi}; this path is selected by
\texttt{loss: m3n\_hinge} together with
\texttt{inference.train: viterbi}, and the configuration
validator rejects mismatched pairings at load time.

\paragraph{LP-relaxed M\textsuperscript{3}N (\texttt{lp\_m3n}).}
The training loss is the LP-relaxed surrogate
$\Delta^{\mathrm{LP}}(x^{i}, y^{i}, w, \phi^{i})$ of
Section~\ref{subsec:lp-m3n}, with each training example carrying
its own auxiliary variables $\phi^{i}$ (stored in the
$\phi$-bank of Section~\ref{subsec:method-training-phi}). No
inference oracle is invoked at training time --- the $\max$
operations are local and evaluated in closed form. This is the
training path used on the Sudoku graph; it is also exercised on
chains in the tightness-validation experiment of
Chapter~\ref{chap:experiments}. The corresponding implementation
is the \texttt{lp\_m3n\_loss} module, selected by
\texttt{loss: lp\_m3n} together with \texttt{inference.train: lp}.

\paragraph{Common scaffolding.}
The remainder of the fit loop is identical for both losses.
Each epoch shuffles the training indices via
\texttt{torch.randperm($N$)}, slices them into mini-batches of
size $B$, and for each mini-batch performs one forward pass
through the backbone, one loss-and-backward step, and one
optimizer step. After each epoch the trainer evaluates Hamming
and $0/1$ on the training and validation splits, records
diagnostics, and decides whether to early-stop.

\subsection{Optimizer and learning-rate schedule}
\label{subsec:method-training-optimizer}

The optimizer is selectable in YAML (\texttt{optimizer.type});
in every experiment of this thesis it is \texttt{adam}, with
SGD supported as an alternative. The optimizer splits the
trainable parameters into up to three groups, each with its own
learning rate:
\begin{itemize}
    \item Group $0$ --- the backbone parameters $\theta$, with
    learning rate $\eta_{0}$ (\texttt{optimizer.lr}) and weight
    decay \texttt{optimizer.weight\_decay}.
    \item Group $1$ --- the pairwise matrix $W$, with learning
    rate \texttt{optimizer.lr\_pairwise} and the same weight
    decay. This group is present for both losses.
    \item Group $2$ --- the per-example duals $\phi^{i}$
    (LP-M3N only), with learning rate \texttt{optimizer.lr\_phi}
    and weight decay \texttt{optimizer.weight\_decay\_phi}, both
    discussed in Section~\ref{subsec:method-training-phi}.
\end{itemize}
The motivation for splitting $W$ off into its own group is the
asymmetric training dynamics between the high-parameter-count
backbone and the small ($K \times K$) pairwise matrix, discussed
empirically in Chapter~\ref{chap:experiments}.

The learning rate decays through a \emph{hyperbolic} schedule
\begin{equation}
    \eta_{t} \;=\; \eta_{0} \cdot
    \frac{\mathrm{offset}}{\mathrm{offset} + t},
    \label{eq:method-lambda-schedule}
\end{equation}
where $t$ counts epochs and \texttt{offset} is a YAML knob.
With \texttt{offset = 100} the multiplier is $1$ at $t = 0$,
$0.5$ at $t = 100$, and $0.1$ at $t = 900$. The choice is
empirically motivated: a fast initial decay locks in the
unary classifier within the first few tens of epochs while the
slow tail preserves gradient signal for the pairwise term,
which is harder to learn.

\subsection{Per-example dual variables}
\label{subsec:method-training-phi}

The LP-M3N loss of Section~\ref{subsec:lp-m3n} requires a set
of per-example auxiliary variables
$\phi^{i} \in \mathbb{R}^{2 \times |E| \times K}$, one set per
training example $i$. The \texttt{Trainer} pre-allocates this
\emph{$\phi$-bank} at construction time (zero-initialized;
persists across epochs), and passes it to the optimizer as a
separate parameter group with its own learning rate
\texttt{lr\_phi}. The motivation for a separate rate is the
asymmetric update frequency under mini-batch SGD: each model
parameter is updated $\sim N / B$ times per epoch, but each
$\phi^{i}$ is updated exactly once (its example appears in
exactly one mini-batch per epoch). Without compensation, a
single global learning rate either inflates $\phi$ in lockstep
with the backbone or starves the dual of signal.

\subsection{Hyperparameter selection}
\label{subsec:method-training-hyperparams}

Hyperparameters (learning rates, weight decays, schedule
offset, batch size, number of epochs) are selected on the
held-out validation split. Each experiment YAML names a single
scalar metric to monitor; training is stopped when this metric
fails to improve for \texttt{patience} consecutive evaluation
ticks, and the model state is reverted to the best-monitored
snapshot on exit. The $\phi$-bank is not snapshotted and
remains at its last-epoch values, which is harmless because
evaluation-time decoding uses only the model's unary and
pairwise outputs.